\subsection{度量与度量空间}

定义 1. 集合 X 上的度量是 X 上的一个有限伪度量 d，使得关系 \(d(x, y) = 0\) 蕴含 x = y。度量空间是配备了由 X 上一个给定度量所定义的结构的集合 X。

度量空间 X 总是被认为带有由 X 上给定度量定义的一致结构和拓扑。

例。1) 欧氏距离 \(d(\pmb{x}, \pmb{y})\) （第六章，§ 2，第 1 号）是实 \(n\) 维数空间 \(\mathbf{R}^n\) 上的一个度量；函数

\[
\sup _ {1 \leqslant i \leqslant n} | x _ {i} - y _ {i} | \quad \text { and } \quad \sum_ {i = 1} ^ {n} | x _ {i} - y _ {i} |.
\]

这些度量都是等价的（§ 1，第 2 号）。

\begin{enumerate}
\def\labelenumi{\arabic{enumi})}
\setcounter{enumi}{1}
\tightlist
\item
  在任意集合 X 上，伪度量 \(d\) ——由关系 \(d(x, x) = 0\) 与 \(d(x, y) = 1\) （当 \(x \neq y\) 时）定义的——是一个度量。它在 X 上定义的一致结构是离散一致结构。
\end{enumerate}

如果我们把度量说成是其定义的一致结构为豪斯多夫的有限伪度量，就得到一个与定义 1 等价的定义；因此，与某个度量等价的有限伪度量是度量。

由单个伪度量（我们可以假定它是有限的）定义的一致空间，当该伪度量不是度量时，可以归化为度量空间。设 \(f\) 是集合 \(\mathbf{X}\) 上的这样一个伪度量，\(\mathcal{U}\) 是由 \(f\) 定义的一致结构；\(\mathcal{U}\) 不是豪斯多夫的，且 \(\mathcal{U}\) 的诸围域之交是 \(\mathbf{X} \times \mathbf{X}\) 的由等价关系 \(f(x,y) = 0\) 所定义的子集。用 \(\mathbf{R}\) 表示这个关系。若 \(x \equiv x' \pmod{\mathbf{R}}\) ，则由三角不等式有

\[
f (x, y) \leqslant f (x, x ^ {\prime}) + f (x ^ {\prime}, y) = f (x ^ {\prime}, y)
\]

同理 \(f(x', y) \leqslant f(x, y)\) ，故 \(f(x, y) = f(x', y)\) ；换句话说，\(f\) 是关于 \(x\) 与 \(y\) 与等价关系 \(\mathbf{R}\) 相容的函数（《集合论》，\(\mathbf{R}\) ，§ 5，第 7 号）。设 \(\overline{f}\) 是 \(f\) 在商集上诱导的函数；\(\overline{f}\) 定义在 \((\mathbf{X}/\mathbf{R}) \times (\mathbf{X}/\mathbf{R})\) 上，且若 \(x\) 与 \(y\) 是 \(\mathbf{X}\) 中任意两点，而 \(\dot{x}\) 与 \(\dot{y}\) 分别记（模 \(\mathbf{R}\) ）之下 \(x\) 与 \(y\) 的等价类，则有 \(\overline{f}(\dot{x}, \dot{y}) = f(x, y)\) 。由此立即可得 \(\overline{f}\) 是 \(\mathbf{X}/\mathbf{R}\) 上的一个度量；它称为与伪度量 \(f\) 相伴的度量；此外，它在 \(\mathbf{X}/\mathbf{R}\) 上定义的一致结构正是按该一致结构的定义与 \(\mathcal{U}\) 相伴的豪斯多夫一致结构（第二章，§ 3，第 8 号，注）。这样，通过过渡到适当的商空间，由单个伪度量定义的一致结构可以归化为度量空间的结构。

第 1 节，第 3 号的命题 1 确定了度量空间完备化的结构：

命题 1. 设 X 是一个度量空间，d 是它的度量。若 \(\hat{X}\) 是 X 的完备化（关于由 \(d\) 定义的一致结构），则函数 \(d\) 可以用连续性扩张到 \(\hat{X} \times \hat{X}\) ；扩张后的函数 \(\overline{d}\) 是 \(\hat{X}\) 上的一个度量，且 \(\hat{X}\) 的一致结构与由度量 \(\overline{d}\) 定义的一致结构重合。

第 1 节，第 3 号的命题 1 表明，\(\overline{d}\) 是 \(\hat{\mathbf{X}}\) 上的有限伪度量，且由 \(\overline{d}\) 在 \(\hat{\mathbf{X}}\) 上定义的一致结构就是由完备化得到的一致结构；由于后一一致结构是豪斯多夫的，故 \(\overline{d}\) 是度量。每当我们把度量空间 \(\mathbf{X}\) 的完备化看作度量空间时，总应理解为 \(\hat{\mathbf{X}}\) 上的度量是由 \(\mathbf{X}\) 上的度量用连续性扩张所得的那个度量。

